\section{Exploratory Analysis Findings}
\label{sec:prototype:eda}

The exploratory data analysis is captured in four executed
Jupyter notebooks~\cite{jupyter} in the project's
\texttt{notebooks/} directory. The notebooks are the empirical
ground truth on which every design decision in
Chapter~\ref{chapter:proposed} rests; each notebook is executed
via \texttt{jupyter nbconvert --execute --inplace} and committed
with its rendered outputs embedded, so report figures can
reference notebook cells directly.

\paragraph{Notebook 01 --- Schema and volume.}
Surveys the parquet exports produced by
\texttt{src/duckdb\_loader.py} and reports the two headline
volume figures that bound the design space: 17{,}883{,}623 AFC
boarding events over the ninety-day window
(2026-01-01 to 2026-03-31), and 720{,}088 GPS pings for the
local one-day sample (2025-03-18). The notebook also verifies
the cross-source overlap (one hundred and twenty-one of one
hundred and sixty-three routes appear in both sources) which
grounds the GPS$\leftrightarrow$AFC reconciliation rule in
\texttt{dim\_route} (Section~\ref{sec:proposed:warehouse_local}).

\paragraph{Notebook 02 --- GPS-side analysis.}
Profiles the local one-day GPS sample. Confirms the
extracted-trips buffer contains nine routes (29, 41, 57, 61, 84,
99, 102, 141, 151), enumerates the one hundred and thirty-one
distinct vehicles, and sanity-checks the door-up and door-down
sensor fields that underpin the ghost-stop and ghost-run
detection planned for the Spatial and Financial cards.

\paragraph{Notebook 03 --- AFC fare and revenue analysis.}
Produces the most consequential finding of the EDA: thirty-four
percent of all boardings on weekday mornings use a student-fare
ticket (6{,}162{,}862 HS-SV tickets out of the 17.8~million
total), and 96.2\% of payments are cash-equivalent. The notebook
identifies the eight universities that anchor the student-fare
share (HCMUS, HCMUT, UEH, GTVT, USSH, HCMUTE, UIT, SGU),
computes the quarterly revenue of approximately
\textcent~72.4~billion, and detects the visible T\d ết dip from
28~January to 5~February 2026.

\paragraph{Notebook 04 --- Stops, schools, and metro.}
Documents the 4{,}710 stops with verified latitude/longitude
and the nineteen metro feeder routes referenced in the dataset
specification~\cite{hcmcgps2025}. Identifies the gap in school
geocoding (the sixty-four school addresses arrive without
coordinates, so \texttt{dim\_school} ships with district-level
stubs initially) and the schedule-data anomaly of ninety-one
thousand \texttt{service\_type=`unknown'} rows that requires
cleaning before any adherence-related chart can render.

\begin{figure}[H]
\centering
\imageplaceholder{Chart (Calendar heatmap)}{A calendar heatmap
of student-fare boardings over the 90-day AFC window, rows =
weeks, columns = days-of-week, colour = student-fare boarding
count. The morning-peak compression on weekdays is the visible
finding; the T\d ết flatline (28~January--5~February) is the
secondary finding. Generated from notebook 03.}{Run
\texttt{notebooks/03\_fare\_revenue.ipynb} cell labelled
``student fare weekly heatmap''. Export the resulting plotly
figure as PNG at $\geq$1200~px wide. Save as
\texttt{images/screenshots/eda\_student\_fare\_calendar.png}.}
\caption{Student-fare boardings by date and weekday, from
Notebook~03. The morning-peak compression on weekdays is the
empirical basis for the flagship card
(Section~\ref{sec:proposed:sdvf}); the T\d ết flatline is the
natural-experiment window for Anomaly Lab Chart~5.5.}
\label{fig:prototype:eda_student_fare}
\end{figure}

The five EDA findings that drove design decisions are documented
in Section~\ref{sec:proposed:eda_to_design}; this section
documents the artefacts (the notebooks themselves) that
\emph{produced} those findings, satisfying the Phase~1
deliverable ``Phát triển, thử nghiệm và đánh giá ban đầu''
required by the capstone topic description~\cite{capstoneDescription}.
